% Vérifié : heure et citations alignées sur la transcription (27 sept. à 11 h ; « tres -- émus », « arret », « Région très boisée ») ; guillemets retirés des formules non littérales (Le Bourget, Invalides) ; « première grande bataille de la Somme » corrigé (Montauban et Trônes Wood en juillet) ; train de 10 h 45 « espéré », non manqué ; décompte des six jours de permission nuancé ; chiffres non vérifiables de la ligne Paris-Toulouse remplacés par le tarif 1912 (R6) ; répétitions de l'introduction générale allégées (études, plaidoirie, paie mensuelle) ; Doullens et nuits en gare reformulés selon R4/R6.
Le 25~septembre 1916, la 152nd Company de l'Army Service Corps bivouaque depuis deux jours en plein air, au flanc d'un coteau près de la route de Forceville, au nord-ouest d'Albert ; Combes s'y est aménagé une hutte (ff.~24~v\textsuperscript{o}--25~r\textsuperscript{o}). La 54th Brigade, qu'elle ravitaille, monte ce jour-là aux abris du bois de Thiepval, et le quartier général de la 18th Division passe d'Acheux à Hédauville : l'assaut de Thiepval est fixé au lendemain\footnote{G.~H.~F.~Nichols, \textit{The 18th Division in the Great War}, Édimbourg et Londres, Blackwood, 1922, p.~79--81 ; Bedfordshire Archives Service, X550/8/1/1609, \textit{War Diary}, 7th (Service) Battalion Bedfordshire Regiment, septembre 1916, 23--25~sept. (cité ci-après JMO 7/Bedf).}. C'est à cette veille de bataille que Combes part en permission ; il ne rejoint sa compagnie que le 6~octobre, dans un autre secteur. Pour la première fois, le carnet quitte la zone des armées et devient, douze jours durant, le journal d'un voyageur et d'un fiancé : deux traversées de Paris, une semaine à Toulouse, ville des études et de la jeunesse de Combes, auprès de Marguerite Calmels. Les comptes du soir, tenus avec le même soin qu'au front, enregistrent désormais des billets de chemin de fer, des tramways, des livres et des bijoux.

\subsection*{Le carnet}

Le matin du départ, Combes doit encore aller chercher du fourrage vert à Forceville avec le lieutenant Ford\footnote{Sans doute Edward Maurice Ford, \textit{2nd Lieutenant} de l'ASC en 1915, qui servit au train de la 18th Division (Nichols, \textit{op.~cit.}, p.~432 ; Spink, \textit{Orders, Decorations and Medals e-Auction~3}, vente 21102, lot~188, spink.com/lot/21102000188).}. À midi, il prend place sur le dernier chariot du convoi qui va à la gare de ravitaillement, gagne à travers champs la gare d'Acheux-Varennes, laisse ses colis au bureau du \textit{Railway Transport Officer} et attend longtemps, « chaleur intense », dans le « petit train départemental » (f.~25~v\textsuperscript{o}). Celui-ci part vers 15~h~30. Combes change de train à Gézaincourt, atteint Amiens vers 20~h et y monte dans un train de permissionnaires pour Paris, en troisième classe, « av grossiers poilus, barbus \& sales ». Au croisement de Longueau, le train est survolé par des « Taubes » ; le ciel est « fouillé par les projecteurs. Très beau spectacle » (f.~26~r\textsuperscript{o})\footnote{« Taube », nom d'un monoplan allemand du début de la guerre, désignait alors tout avion allemand : E.~Fraser et J.~Gibbons, \textit{Soldier and Sailor Words and Phrases}, Londres, Routledge, 1925, s.v. Ce raid n'a été retrouvé dans aucune des sources consultées.}. Vers 6~h, il doit descendre en gare du Bourget : il suit les voies jusqu'au bureau militaire, y obtient l'autorisation de sortir et un planton l'accompagne jusqu'à la porte (f.~26~v\textsuperscript{o}).

La journée parisienne du 26~septembre est faite de démarches et de flânerie. De la gare d'Orsay, Combes se rend à l'hôtel des Invalides pour demander l'autorisation de voyager « à mes frais en express », avant de prendre un billet de deuxième classe pour Toulouse, payé 15~F au « 1/4 place »\footnote{Le gouverneur militaire de Paris siégeait aux Invalides : K.~Baedeker, \textit{Paris and Environs}, 18\textsuperscript{e}~éd., Leipzig, Baedeker, 1913, p.~263--264. Le texte qui réglait le passage des permissionnaires dans les express n'a pas été retrouvé.}. Il passe ensuite au Bon Marché, au Luxembourg et « Chz Colin », se repose au jardin de Notre-Dame et s'attarde sur un banc du square de la tour Saint-Jacques : « grd mouv\textsuperscript{t} rue de Rivoli. Animation ordin\textsuperscript{re} » (ff.~26~v\textsuperscript{o}--27~r\textsuperscript{o}). Les comptes de la journée en font l'inventaire : un képi au Bon Marché (8,75~F), un objet de chez Christofle (10,50~F après remise), une couronne de fleurs artificielles aux magasins de l'Hôtel-de-Ville, une carte de Verdun et une carte de Roumanie, et deux livres, \textit{L'Orgueil humain} de Zyromski et \textit{Les Débris de la guerre} de Maeterlinck (f.~27~r\textsuperscript{o})\footnote{E.~Zyromski, \textit{L'Orgueil humain}, Paris, A.~Colin, 1904 ; M.~Maeterlinck, \textit{Les Débris de la guerre}, Paris, Fasquelle, 1916, porté à 3,50~F, le prix noté par Combes, dans les listes de l'éditeur (exemplaire de 1918, archive.org/details/lesdbrisdelagu00maetuoft). «~Colin~» est vraisemblablement la librairie Armand Colin, 103, boulevard Saint-Michel, adresse imprimée au titre de ses publications.}. Le train quitte Orsay vers 20~h et arrive à Toulouse le 27 à 7~h~30\footnote{Le plein tarif de Paris à Toulouse était alors de 54,30~F en deuxième classe : K.~Baedeker, \textit{Le Sud-Ouest de la France, de la Loire à la frontière d'Espagne}, Leipzig, Baedeker, 1912, route~39.}.

À Toulouse, Combes retrouve son appartement --- celui de la rue de Languedoc, dont il paie le 3~octobre à sa logeuse, Mlle Guilhem, le semestre de juillet 1916 à janvier 1917 (90~F, f.~31~r\textsuperscript{o})\footnote{Le carnet ne nomme pas la rue pendant ce séjour ; elle l'est le 1\textsuperscript{er}~février 1917 (f.~67~v\textsuperscript{o}) et, avec le n\textsuperscript{o}~24, dans une carte-lettre de Combes à Marguerite, Paris, 25~novembre 1918, archives familiales.}. Plusieurs matinées commencent par la même scène : il s'attarde chez lui, prend « un plaisir extrême à regarder \& ranger [ses] affaires » et laisse « s'éveiller [...] des souvenirs du passé » (f.~28~r\textsuperscript{o} ; de même ff.~29~v\textsuperscript{o} et 30~r\textsuperscript{o}). Le 27, vers 11~h, il « monte » à Sainte-Marie-des-Champs, sur la Côte-Pavée, chez les Crouzet, oncle et tante de Marguerite ; elle et sa mère n'y sont attendues que pour le soir, mais il la revoit dès 17~h, « Charmante de simplicité \& de tendresse confiante » (f.~27~v\textsuperscript{o}). Les après-midi sont faits de promenades, en tramway puis à pied : Saint-Agne, Pouvourville et la crête de Pech-David le 28 ; Purpan, ses « Ruines » et Blagnac le 29 ; le clocher de Montaudran, d'où il fait « vibrer chacune des 5 cloches d'airain », puis le cimetière de Terre-Cabade le 30 ; Lardenne et les bords du Touch le 1\textsuperscript{er}~octobre (ff.~28~r\textsuperscript{o}--29~v\textsuperscript{o}). La géographie religieuse de ces journées est dense : messe à Saint-Étienne le 1\textsuperscript{er}~octobre, « M\textsuperscript{te}, à ma droite » ; chapelle de Nazareth le 2 ; le 3, à la cathédrale, messe anniversaire pour une « Mlle Calmels » que les sources n'identifient pas (ff.~29~v\textsuperscript{o}--30~v\textsuperscript{o}). À Terre-Cabade, devant le caveau des jésuites et celui du père Marie-Antoine, les fiancés échangent des serments « au souvenir de notre gd ami \& de Josephe », c'est-à-dire de Joséphine, la sœur de Combes morte en 1914 (f.~29~r\textsuperscript{o}) ; l'introduction générale a commenté l'épisode.

Autour du couple, le carnet laisse apparaître un cercle. Il y a les Crouzet, chez qui il déjeune le premier jour « ds la + gd intimité » (f.~27~v\textsuperscript{o}), « Jean », qui le raccompagne le soir en tramway les 1\textsuperscript{er} et 2~octobre, et Mme de Labarthe, « gde amie affectueuse \& attachée », qu'il visite deux matins de suite et à qui il présente « ma Fiancée » le 3, dans une propriété dont elle leur fait visiter la maison (ff.~28~v\textsuperscript{o}--31~r\textsuperscript{o})\footnote{« Jean » est sans doute l'abbé Jean Crouzet (v.~1880--1932), vraisemblablement fils d'Émile Crouzet, qu'on trouve vicaire à Saint-Étienne pendant la guerre : \textit{Le Télégramme} (Toulouse), 5~février 1915, p.~2, et 15~avril 1917, p.~3 (Rosalis, Bibliothèque de Toulouse). Le terminus de la promenade du 3~octobre («~Castanet~» ?) est d'une lecture incertaine.}. Le 2, Marguerite et sa mère viennent le voir chez lui ; le soir, « Ts 3 », ils regardent les vitrines de la rue d'Alsace (f.~30~r\textsuperscript{o}). Le 4, jour du départ, il range l'appartement, « Enlevé \& rentré rideaux », passe l'après-midi à Sainte-Marie, « Joué avec des enfants » dans le parc. Marguerite veut « absol\textsuperscript{t} » l'accompagner à la gare avec l'oncle Crouzet, et il faut se séparer « par ménagement pour l'oncle » ; « M\textsuperscript{t} \& moi, tres -- émus » (ff.~31~v\textsuperscript{o}--32~r\textsuperscript{o}).

Le retour est plus pénible que l'aller (\textit{Feuillets de Route}, chap.~\textsc{xxii} : «~Retour de permission au front~», annexe, p.~\pageref{annexe:doc6}). Combes a pris pour Paris un billet de troisième classe (11,75~F, f.~32~r\textsuperscript{o}) et ne parvient à s'assoupir qu'au matin. À Orsay le 5 à 8~h, il traverse Paris sous une pluie fine, laisse ses paquets à la consigne de la gare du Nord, admire à Saint-Vincent-de-Paul les « peintures murales et rangées superposées de belles colonnes de marbre marron », écrit deux lettres au salon de correspondance des Grands Magasins du Louvre et, « Trop long arret au Louvre », n'a plus le temps de déjeuner (f.~32~v\textsuperscript{o})\footnote{L'église, commencée en 1824 par J.-B.~Lepère et achevée par J.-I.~Hittorff, est connue pour ses deux étages de colonnades et la frise peinte d'Hippolyte Flandrin (1850--1854) : base Mérimée, notice PA00086489 ; Baedeker, \textit{Paris and Environs}, \textit{op.~cit.}, p.~51--52.}. Il avait espéré un train à 10~h~45 ; il part à 13~h~15, est à Amiens à 18~h et atteint Doullens vers minuit : « Dû coucher gare » (f.~33~r\textsuperscript{o}). Il achève la nuit « allongé sur 1 banc » de la salle des départs, erre au petit jour dans la ville encore déserte et entre dans l'église pour y rêver « à [son] aise ». Au train de 7~h~25, le R.T.O. lui apprend « que ma division a shifté. Le railhead est Canaples ». Il y trouve les camions de la \textit{supply column} de la 54th Brigade, monte dans l'un d'eux jusqu'au \textit{refilling point} et retrouve la 152nd Company à 13~h (f.~33~v\textsuperscript{o}).

\subsection*{Le contexte}

Pendant ces douze jours, la division de Combes livre la bataille de Thiepval. Le 25~septembre, jour de son départ, la Fourth Army prend Morval et Lesbœufs\footnote{W.~Miles, \textit{Military Operations, France and Belgium, 1916}, vol.~II, Londres, Macmillan, 1938, p.~373 (cité ci-après \textit{OH 1916}, II).}. Le 26, à 12~h~35, les 53rd et 54th Brigades de la 18th Division attaquent Thiepval ; le village est entièrement pris à l'aube du 27 par deux compagnies du 7th Bedfordshire. La division dispute ensuite la Schwaben Redoubt jusqu'au 5~octobre, puis la 39th Division la relève. Du 26~septembre au 5~octobre, elle a perdu 3\,344~hommes\footnote{Nichols, \textit{op.~cit.}, p.~79--82 ; \textit{OH 1916}, II, p.~406--417 et 453, n.~5 ; JMO 7/Bedf, 26--30~sept. 1916 ; A.~F.~Becke, \textit{Order of Battle of Divisions}, part~3A, Londres, HMSO, 1938, p.~84.}. Le 2~octobre, le 7th Bedfordshire est transporté en wagons de la gare d'Acheux à Candas, et la division se repose dans la région de Candas, à l'ouest de Doullens\footnote{JMO 7/Bedf, octobre 1916 (X550/8/1/1610), 2~oct. ; Nichols, \textit{op.~cit.}, p.~126.}. C'est ce déplacement que Combes apprend à Doullens ; de Paris ou de Toulouse, il n'a mentionné ni Thiepval ni les pertes de sa division. Les cartes de Verdun et de Roumanie achetées le 26 attestent seulement qu'il suit les nouvelles des autres fronts. La Roumanie était entrée en guerre un mois plus tôt\footnote{Déclaration de guerre à l'Autriche-Hongrie du 27~août 1916 : \textit{La Dépêche} (Toulouse), 29~août 1916, p.~1.}.

Le voyage suit le circuit du permissionnaire, alors en pleine expansion. Les trains spéciaux de permissionnaires circulaient depuis l'automne 1915 ; le réseau du Nord en fit rouler 4\,200 en 1915, 11\,800 en 1916. Le Bourget et Creil, gares régulatrices des armées françaises depuis février 1915, venaient d'être agrandies pour l'offensive de la Somme, et le service des permissionnaires y avait exigé « de vastes et nombreuses installations » : de là le contrôle et l'escorte du 26 au matin\footnote{G.~Lafon, \textit{Les Chemins de fer français pendant la guerre}, thèse de droit, Paris, A.~Rousseau, 1922, p.~99, 105, 115 et 152--153.}. Le « quart de place » applique l'article~54 des cahiers des charges des compagnies, qui réduisait au « quart de la taxe » le prix payé par les militaires en permission\footnote{\textit{Ibid.}, p.~37.}. Le 3~novembre, la Mission retiendra à Combes « 30~f [...] pour 6 jours de permission » (f.~46~r\textsuperscript{o}), soit 5~F par jour, vraisemblablement sur son indemnité de fonctions, puisque le prêt du soldat était maintenu pendant la permission. Ce décompte est celui de la règle alors en vigueur, « six jours pleins », trajet non compris\footnote{\textit{Bulletin des armées de la République}, 2\textsuperscript{e}~année, n\textsuperscript{o}~186, 22~mars 1916, p.~2, «~Les permissions~». Le régime de trois permissions annuelles de sept jours ne s'appliqua qu'à partir du 1\textsuperscript{er}~octobre 1916 : E.~Cronier, \textit{Permissionnaires dans la Grande Guerre}, Paris, Belin, 2013, p.~18 et 32, cité d'après l'article «~Permissionnaires français de la Première Guerre mondiale~» de Wikipédia.}, bien que le séjour à Toulouse s'étende du 27~septembre au 4~octobre. Au retour, il fait viser son titre de permission au bureau de la Mission à Bernaville (f.~34~r\textsuperscript{o}).

Les comptes de Toulouse, faits surtout de centimes --- tramways, bains, fruits, fleurs, images pieuses ---, comportent trois dépenses exceptionnelles : un sautoir en or chez le bijoutier Pfister, « net 295~f » (f.~28~v\textsuperscript{o}), la façon d'une tunique chez le tailleur Guiraud (38~F) et le loyer du semestre (f.~31~r\textsuperscript{o}). Les versements de la Mission font 157,50~F en septembre 1916 ; le sautoir représente près de deux mois de cette paie, la tunique environ quatre journées et demie d'un tailleur d'habits parisien, payé 8,50~F par jour en 1916\footnote{Totaux mensuels calculés d'après les versements notés au carnet ; L.~March, \textit{Mouvement des prix et des salaires pendant la guerre}, Paris, PUF, et New Haven, Yale University Press, 1925, p.~264.}. Ces achats supposent une épargne ou d'autres ressources, que le carnet ne fait pas connaître. Les tarifs du métro, eux, sont restés ceux de 1913 : 15~c en deuxième classe (f.~33~r\textsuperscript{o})\footnote{Baedeker, \textit{Paris and Environs}, \textit{op.~cit.}, p.~28--30.}.

Toulouse comptait 149\,576~habitants en 1911 et était le siège du XVII\textsuperscript{e} corps d'armée, d'un archevêché et d'une université\footnote{Baedeker, \textit{Le Sud-Ouest de la France}, \textit{op.~cit.}, p.~298--299.}. De la ville de guerre, le carnet ne retient rien ; Combes y circule dans un quartier familier. Les Thermes du Languedoc, où il prend ses bains, sont au n\textsuperscript{o}~1~bis de sa rue, et Pfister tient boutique au 24, place des Carmes, « angle rue Languedoc »\footnote{\textit{Le Télégramme}, 31~mars 1911, p.~4 ; \textit{L'Express du Midi}, 25~décembre 1911, p.~3 (Rosalis).}. Les tramways électriques de la Compagnie Pons le mènent aux terminus de la banlieue, place Esquirol pour Saint-Agne, Purpan et Lardenne, Capitole pour la Côte-Pavée\footnote{Baedeker, \textit{Le Sud-Ouest de la France}, \textit{op.~cit.}, p.~298 ; Wikipédia, «~Ancien tramway de Toulouse~», d'après C.~Lacombe et P.~Malterre, \textit{Omnibus, tramways et autobus de Toulouse}, Toulouse, éd. du Cabri, 1983.}. Sainte-Marie-des-Champs portait le nom de l'ancien noviciat des jésuites, établi vers 1843 dans une maison de campagne « à une demi-heure de la ville ». Une « villa Sainte-Marie-des-Champs », au terminus du tramway de la Côte-Pavée, apparaît encore dans la presse locale de 1914 à 1925 ; la maison des Crouzet doit probablement y être située\footnote{J.~Burnichon, \textit{La Compagnie de Jésus en France. Histoire d'un siècle, 1814--1914}, t.~II, Paris, Beauchesne, 1916, p.~297--298 ; \textit{La Dépêche}, 27~juin 1914, p.~4 ; \textit{Le Midi socialiste}, 5~décembre 1922, p.~3. Le statut des Crouzet dans les lieux n'est pas connu.}. Le capucin Marie-Antoine, figure populaire du catholicisme toulousain, avait été inhumé à Terre-Cabade en février 1907 devant une foule de « plus de cinq mille personnes »\footnote{\textit{L'Express du Midi}, 11~février 1907, p.~2.}. La chapelle de Nazareth, rue Philippe-Féral, accueillait des messes anniversaires et des services pour des morts de la guerre\footnote{\textit{Le Télégramme}, 10~mars 1916, p.~4, et 3~juillet 1917, p.~2.}.

La nuit à Doullens est la première des nuits en gare qui marqueront les voyages de permission. La ville, dont la citadelle abritait un hôpital militaire, était depuis 1916 l'un des centres sanitaires de l'arrière britannique, où s'installèrent plusieurs \textit{Casualty Clearing Stations}\footnote{Commonwealth War Graves Commission, notice «~Doullens Communal Cemetery Extension No.~1~», cwgc.org.}. Le trajet du 6~octobre, de la gare de ravitaillement de Canaples au \textit{refilling point} de la brigade, ramène Combes par la chaîne même du ravitaillement qu'il sert depuis l'été, et ouvre le chapitre suivant. La compagnie cantonne désormais à Barlette, près de Ribeaucourt, « Région très boisée, ext\textsuperscript{t} pittoresque », où Combes loge chez une veuve (f.~34~r\textsuperscript{o}). Ce n'est qu'une halte : la division se reconstitue après Thiepval avant de revenir sur l'Ancre à la mi-octobre\footnote{JMO 7/Bedf, octobre 1916 ; Nichols, \textit{op.~cit.}, p.~126 et 130.}.
